\einheitenkopf{Einheit 1 von 5 · Häufigkeiten}
\erg{8}{a) 18 \quad b) 250 \quad c) 26 \quad d) 40 \quad e) 25}
\erg{9}{a) 8 \quad b) 12 \quad c) 4 \quad d) 16 \quad e) Anzahlen: 1 → 3, 2 → 2, 3 → 7, 4 → 2, 5 → 3, 6 → 3 (zusammen 20)}
\erg{10}{a) $\tfrac{3}{10}$ \quad b) $\tfrac{7}{20}$ \quad c) $\tfrac{9}{25}$ \quad d) $\tfrac{4}{50}$ \quad e) $\tfrac{6}{24}=\tfrac14$ \quad f) $0{,}65$ \quad g) $0{,}425 = 42{,}5\,\%$ \quad h) $5:9 = 0{,}5555\ldots \approx 55{,}6\,\%$ \quad i) $45+35=80$, fehlender Anteil $20\,\%$ \quad j) 9 von 40 zusammen $\rightarrow 22{,}5\,\%$ \quad k) $0{,}15 \cdot 320 = 48$ \quad l) $h=\tfrac{26}{40}=\tfrac{13}{20}=0{,}65=65\,\%$}
\erg{11}{a) Mia hat die Anzahl hingeschrieben, nicht den Anteil: 8 ist die absolute Häufigkeit. Richtig ist $h = 8 : 32 = 0{,}25 = 25\,\%$. \quad b) $\tfrac{18}{45}=\tfrac25 = 0{,}4 = 40\,\%$}
\erg{12}{a) Die Anzahlen hängen von der Gruppengröße ab; erst das Teilen durch die Gesamtzahl macht die Gruppen vergleichbar. \quad b) In beiden Klassen ist der Anteil gleich groß: jeweils die Hälfte, also $50\,\%$.}
\erg{13}{a) Nordschule $84:240 = 0{,}35 = 35\,\%$, Südschule $78:195 = 0{,}4 = 40\,\%$ \quad b) Nordschule (84 gegen 78) \quad c) Südschule \quad d) Südschule; die Anzahl allein entscheidet nicht, weil die Nordschule viel mehr Kinder hat.}
